% ------------------------------------------------------------- PX4 SITL transfer study
\subsection{RQ7: transfer to PX4}
\label{sec:px4}
\paragraph{Harness.}
We repeated the key experiments in PX4 v1.14.3 SITL with its SIH physics, flying
a 60\,s straight segment at about 8\,m/s in mission mode, \nPxMissions{} missions in all.
The GNSS offset and the benign envelope (gust, GNSS noise, outages, jamming,
IMU faults, LoE) are injected inside the simulator from seeded streams, so each
attacked mission has a twin with the same noise and the attack off; the
attack-only deviation is the largest difference between the two. Stock EKF2 v1.14
fuses vision position only relative to a bias estimate while GNSS is active, so
a trusted fix would be absorbed; a small EKF2 change fuses the fix every 5\,s as
an absolute position. D1--D4 and the anchor gate are computed offline from the
flight logs, with D2 taken from EKF2's own GNSS innovation test, and thresholds
are recalibrated on \nPxBenign{} benign missions. The protocol was pre-registered
in the artifact before the first attacked flight. Two dated deviations
followed. In \nPxNotFlownMin--\nPxNotFlownMax{} of missions a start-up race in
the harness sent PX4 into a failsafe before the window, independently of the
attack, and these missions are excluded. The PX4-optimised attacker of each family
was re-selected after this exclusion. Flying one benign mission twice differs by
\nPxReplayMed\,m in the median and \nPxReplayMax\,m at most, a floor under every
attack-only number below.

\begin{table}[t]
\centering\footnotesize
\caption{PX4 SITL, 5\,s re-anchoring, 5\% FPR. Deviations (m) are attack-only on
undetected missions. RO: reduced-order study (detection: D2).}
\label{tab:px4}
\setlength{\tabcolsep}{2.5pt}
\begin{tabular}{@{}lrr@{}}
\toprule
 & PX4 & RO \\
\midrule
Step 5\,m caught: EKF2 / stack$+$gate & \nPxEOneStepDTwo{} / \nPxEOneStepAll & \nTprStepDTwo \\
Ramp 0.5\,m/s caught: EKF2 / stack$+$gate & \nPxEOneRampDTwo{} / \nPxEOneRampAll & \nTprRampDTwo \\
Transferred attackers: evasion & \nPxEvMin--\nPxEvMax & \nRoPrimEvMin--\nRoPrimEvMax \\
Transferred attackers: largest deviation & \nPxPrimMaxMin--\nPxPrimMaxMax & \nRoPrimMaxMin--\nRoPrimMaxMax \\
Worst of 35, stack$+$gate & \nPxWorstAllGate & \nAnchorMaxFive \\
Worst of 35, no gate: stack / D4 / D1 & \nPxWorstAllNoGate/\nPxWorstDFourNoGate/\nPxWorstDOneNoGate & \\
PX4-optimised: evasion; largest dev. & \nPxCemEvMin--\nPxCemEvMax; \nPxCemDevMax & \\
Stealth budget $\dadv$ & \nPxDadv & \nPxDadvRoPilot/\nPxDadvRoPooled \\
No anchor, ramp: evasion; drift & \nPxNoAnchorEv; \nPxNoAnchorDev & \\
Re-anchoring 1\,s: worst, stack$+$gate / D1$+$gate & \nPxTaWorstAllGate/\nPxTaWorstDOneGate & bound \nBoundTaOne \\
\bottomrule
\end{tabular}
\end{table}

\paragraph{Results.}
\Cref{tab:px4} summarises the study. Against naive offsets, EKF2's own gate
caught the step in \nPxEOneStepDTwo{} and the ramp in \nPxEOneRampDTwo{} of
missions, and the full stack with the anchor gate caught both every time. The
35 attackers optimised in the reduced-order model, flown without
re-optimisation, still evaded the full PX4 stack: the five primary ones in
\nPxEvMin--\nPxEvMax{} of missions, but their median attack-only deviation was
only \nPxQfMin--\nPxQfMax\,m, near the replay floor, and the worst over all
attackers was \nPxWorstAllGate\,m. Attackers optimised against PX4 itself evaded
in \nPxCemEvMin--\nPxCemEvMax{} of missions with at most \nPxCemDevMax\,m; no GNSS
or weather mimic evaded at all. The PX4 stealth budget, \nPxDadv\,m (median
\nPxDadvMed, 95th percentile \nPxDadvQnf\,m), lies between the reduced-order
\nPxDadvRoPilot{} and \nPxDadvRoPooled\,m. Without the trusted fix, a 0.5\,m/s ramp
evaded in \nPxNoAnchorEv{} of missions and moved the vehicle by the full
\nPxNoAnchorDev\,m, as the finiteness result predicts. For the stack of the
proved bound (D1 and the anchor gate), the worst PX4 case was
\nPxWorstDOneGate\,m, inside the \nBoundTaFive\,m bound at 5\,s; this is an
empirical consistency check, since the proof covers the linear loop only. D3's
threshold, set by benign outages, made it as likely to fire on a naive attack as
on a benign mission (\nPxEOneStepDThree{} and \nPxEOneRampDThree).
With re-anchoring every 1\,s (pre-registered with two predictions), both held: the
worst undetected deviation was \nPxTaWorstDOneGate\,m for the bound's stack, inside
the \nBoundTaOne\,m bound and below the 5\,s value, and \nPxTaWorstAllGate\,m for
the full stack with the gate (\nPxTaWorstAllNoGate\,m even without it); the naive
ramp evaded in \nPxTaRampEv{} of missions.

\paragraph{The pre-registered reading did not hold.}
The protocol fixed two conditions for a qualitative transfer of the disguise
matrix. (a)~With the gate armed, every primary attacker should evade in at least
half of the missions with at most twice its reduced-order deviation; \nPxReadAOk{}
of five families met it, but the step reached \nPxReadStepDev\,m against a limit of
\nPxReadStepLim\,m. (b)~With the gate removed, EKF2's gate should cap every family
at \nPxReadCap\,m while D4 does not; D4 admitted \nPxReadFaultDFour\,m to the fault
mimic, but EKF2's gate also admitted \nPxReadFaultDTwo\,m. What transferred is
the ordering: the anchor gate decides how far an undetected attacker gets, the
inertial invariant D4 is bypassed by the fault mimic, and without the gate the
damage is two to three times larger in PX4 than in the reduced-order model.
